\chapter{Quotient кривизны длины два и каноническая высота связывающей цепи}
\label{ch:v7-linking-chain-degree-two-curvature-quotient-gate}

\section{Две разные операции, которые нельзя смешивать}

Трёхступенная суперсвязность предыдущего гейта дала точную факторизацию
\begin{equation}
 (d_A^2)_{20}=B_1B_0=\mathcal R_U,
 \label{eq:v7-degree-two-endpoint-block}
\end{equation}
но её полный самосопряжённый квадрат содержит также диагональные
backtracking-блоки. Теперь необходимо различить:
\begin{enumerate}
 \item стандартный представленный фактор двухформ Конна;
 \item относительное чтение кривизны по степени цепи.
\end{enumerate}
Первое является quotient универсального дифференциального исчисления по
junk-идеалу. Второе является производной mapping-cone типа. Совпадение их
результатов нельзя предполагать.

\section{Обычный junk-фактор удаляет нужный путь}

Рассмотрим узловую алгебру
\begin{equation}
 \mathcal A_{\rm node}=\mathbb C^3
 \label{eq:v7-degree-two-node-algebra}
\end{equation}
на блоках
\begin{equation}
 H_0=\mathbb C^{11},\qquad
 H_1=\mathbb C^{21},\qquad
 H_2=\mathbb C^{10}.
 \label{eq:v7-degree-two-node-blocks}
\end{equation}
Её минимальные идемпотенты действуют как единицы соответствующих блоков, а
оператор Дирака равен $Q_A=d_A+d_A^*$.

Представленные формы вычисляются стандартно:
\begin{align}
 \pi_D(e_i\,de_j)&=e_i[Q_A,e_j],\nonumber\\
 \pi_D(e_i\,de_j\,de_k)&=e_i[Q_A,e_j][Q_A,e_k],
 \label{eq:v7-degree-two-represented-forms}
\end{align}
а junk второй степени равен
\begin{equation}
 J_D^2=\pi_D\!\left(d\ker\pi_D|_{\Omega^1}\right).
 \label{eq:v7-degree-two-junk}
\end{equation}

На двадцати общих комплексных полях $A$ машинный расчёт стабильно даёт
\begin{equation}
 \rank\pi_D(\Omega^1)=4,\qquad
 \rank\pi_D(\Omega^2)=6,\qquad
 \rank J_D^2=2,
 \label{eq:v7-degree-two-ranks}
\end{equation}
и потому
\begin{equation}
 \rank\Omega_D^2=4.
 \label{eq:v7-degree-two-quotient-rank}
\end{equation}
Но оба крайних блока $0\to2$ и $2\to0$ целиком лежат в junk. Максимальный
остаток их проекции вне $J_D^2$ меньше $10^{-10}$.

\begin{nogocriterion}[No-go обычного junk-объяснения]
Стандартный фактор $\Omega_D^2/J_D^2$ не сохраняет
$B_1B_0=\mathcal R_U$, а удаляет его. Поэтому правильный локальный
селектор нельзя объявлять обычной двухформой Конна данного трёхузлового
исчисления.
\end{nogocriterion}

\section{Канонический оператор степени цепи}

Трёхступенный носитель имеет собственный целочисленный оператор степени
\begin{equation}
 N=\operatorname{diag}(0\cdot I_{11},1\cdot I_{21},2\cdot I_{10}).
 \label{eq:v7-degree-two-number-operator}
\end{equation}
Его последовательные разности равны единице. Поэтому $N$ определён
структурой цепи с точностью до добавления скаляра и смены ориентации.
Обе неоднозначности исчезают в норме коммутатора.

Для полной чётной кривизны
\begin{equation}
 F_A=Q_A^2-Q_{A_0}^2
 \label{eq:v7-degree-two-full-curvature}
\end{equation}
определим нормированную относительную производную
\begin{equation}
 \delta_N(F_A)=\frac12[N,F_A].
 \label{eq:v7-degree-two-relative-derivative}
\end{equation}
Диагональные блоки имеют разность степеней ноль и исчезают. Крайние блоки
имеют разность степеней два и сохраняются с единичным коэффициентом:
\begin{equation}
 \delta_N(F_A)=
 \begin{pmatrix}
 0&0&-\mathcal R_U^*\\
 0&0&0\\
 \mathcal R_U&0&0
 \end{pmatrix}.
 \label{eq:v7-degree-two-relative-blocks}
\end{equation}
Следовательно,
\begin{equation}
 \|\delta_N(F_A)\|_{\rm HS}^2
 =2\|\mathcal R_U(A)\|_{\rm HS}^2.
 \label{eq:v7-degree-two-relative-norm}
\end{equation}
На случайных полях максимальные остатки блочного тождества и нормы меньше
$10^{-8}$.

\section{Ковариантность и ориентация}

Оператор $N$ скалярен на каждом узловом блоке и коммутирует с
блок-диагональной gauge-группой. Поэтому
\begin{equation}
 \delta_N(gF_Ag^*)=g\,\delta_N(F_A)\,g^*.
 \label{eq:v7-degree-two-covariance}
\end{equation}
Максимальный машинный остаток ковариантности меньше $10^{-10}$.

Смена ориентации цепи заменяет $N$ на $2I-N$ и меняет знак
$\delta_N$, но не её норму:
\begin{equation}
 \|\delta_{2I-N}(F_A)\|_{\rm HS}
 =\|\delta_N(F_A)\|_{\rm HS}.
 \label{eq:v7-degree-two-orientation-independence}
\end{equation}
Таким образом, выбор физического направления цепи не вводит непрерывного
веса.

\section{Относительное Hodge-действие}

Естественный относительный функционал имеет вид
\begin{equation}
 \mathcal S_{\rm rel}
 =\mathcal S_E+\frac14
 \|\delta_N(Q_A^2-Q_{A_0}^2)\|_{\rm HS}^2
 =\mathcal S_E+\frac12\|\mathcal R_U(A)\|_{\rm HS}^2.
 \label{eq:v7-degree-two-relative-action}
\end{equation}
Коэффициент $1/2$ справа является удалением двух Real-ориентаций крайнего
блока, а множитель $1/2$ в $\delta_N$ фиксирован разностью степеней два.
Непрерывный параметр отсутствует.

Повторный гессиан даёт
\begin{align}
 (n_-,n_0,n_+)_{0}&=(7,0,20),
 &\lambda_{\rm heavy}^{\min}&=\frac{18}{5},
 \label{eq:v7-degree-two-origin-signature}\\
 (n_-,n_0,n_+)_{A_0}&=(0,0,27),
 &\lambda_{\min}&=3.9368554658\ldots.
 \label{eq:v7-degree-two-vacuum-signature}
\end{align}

\begin{theorem}[Канонический relative mapping-cone quotient]
Оператор степени трёхступенной цепи канонически задаёт производную
$\delta_N=\frac12[N,\cdot]$, которая удаляет все диагональные
backtracking-блоки полного квадрата и сохраняет ровно относительную
полярную кривизну $\mathcal R_U$. Полученная положительная норма
калибровочно ковариантна, не зависит от ориентации цепи и воспроизводит
правильные нулевой и вакуумный гессианы без непрерывного секторного веса.
\end{theorem}

\begin{remark}
Это положительный результат для relative mapping-cone наблюдаемого, но не
для обычного Connes-junk quotient: последний удаляет нужный endpoint-класс.
Остаётся вывести рёберный Hodge-момент и $\delta_N(F_A)$ как ортогональные
компоненты одной общей Hodge-метрики, а не как сумму двух уже выбранных норм.
\end{remark}

\begin{nogocriterion}[Запрет переименования mapping cone в junk]
Нельзя ссылаться на обычный junk-идеал как на источник
\eqref{eq:v7-degree-two-relative-action}: его прямой аудит даёт
противоположный результат. Физическое право использовать
$\delta_N(F_A)$ должно следовать из цепной степени, граничного observable
или общей относительной Hodge-структуры.
\end{nogocriterion}

\begin{protocol}
\begin{enumerate}
 \item узловая алгебра: \textbf{$\mathbb C^3$};
 \item ранги $\Omega^1,\Omega^2,J^2$: \textbf{$4,6,2$};
 \item ранг обычного quotient: \textbf{$4$};
 \item endpoint-путь в обычном junk: \textbf{да};
 \item обычный junk-маршрут: \textbf{закрыт};
 \item оператор степени $N=(0,1,2)$: \textbf{каноничен};
 \item $\delta_N$ удаляет возвраты: \textbf{точно};
 \item $\delta_N$ сохраняет $\mathcal R_U$: \textbf{точно};
 \item gauge-ковариантность: \textbf{получена};
 \item ориентационная независимость нормы: \textbf{получена};
 \item нулевая сигнатура: \textbf{$(7,0,20)$};
 \item вакуумная сигнатура: \textbf{$(0,0,27)$};
 \item общий Hodge-след двух компонент: \textbf{открыт};
 \item следующий гейт: \textbf{общий цепно-Hodge relative trace}.
\end{enumerate}
\end{protocol}

Машинный сертификат сохранён в
\path{s2t_v7_linking_chain_degree_two_curvature_quotient_gate_results.json}.

\begin{thebibliography}{9}
\bibitem{v7-degree-two-schucker}
T.~Sch\"ucker, J.-M.~Zylinski,
\emph{Connes' Model Building Kit},
arXiv:hep-th/9312186.
\bibitem{v7-degree-two-roepstorff}
G.~Roepstorff,
\emph{Superconnections and the Higgs Field},
arXiv:hep-th/9801040.
\bibitem{v7-degree-two-project}
Том~IV, гейт junk no-go и канонического mapping-cone селектора.
\end{thebibliography}